\chapter*{Preface}
\addcontentsline{toc}{chapter}{Preface}

\section*{Who this book is for}

This book is for adults who have never had a music lesson. You do not need to know what a note, a chord or a beat is. The book starts from the first key you press and explains every word the first time it appears.

By the end you will know your way around a piano keyboard. You will be able to build scales and chords, keep a steady beat, read simple sheet music, and play your first short pieces with both hands.

\section*{How to use this book}

Music is something you hear, so this book is meant to be heard. In the web edition, every example has a small round play button next to it. Tap it and you hear exactly what the text describes. Try it now with the note called middle C \listen{C4:2}.

The web edition also has keyboards you can tap. When the text says ``press this key'', you can press it on the screen and hear it.

A real instrument helps but is not required. A cheap electronic keyboard or a free piano app on your phone is enough. If you have nothing, read with the play buttons and point at the pictures.

Go slowly. Each chapter is short and builds on the one before. When a box says ``Try it at the keyboard'', give it a minute before reading on.

\section*{A two-minute experiment}

Try this now, even before Chapter~1, which explains the names.

\begin{enumerate}
  \item Find a pair of black keys near the middle of the keyboard. The white key directly to the left of the pair is called C. Put your right thumb on it.
  \item Count four keys to the right, counting every key, black and white, but not the one you start on. Put your middle finger there. That key is E.
  \item Count three more keys to the right. Put your little finger there. That key is G.
  \item Press all three keys together. This is a C major chord. It sounds bright and settled \listen{C4+E4+G4:2}.
  \item Keep your thumb and little finger where they are. Move only your middle finger one key to the left, onto a black key. Press the three keys again. This is a C minor chord. It sounds darker and more thoughtful \listen{C4+Eb4+G4:2}.
\end{enumerate}

Hear the two one after the other \listen{C4+E4+G4:2 C4+Eb4+G4:2}. One finger moved by one key, and the whole mood changed. The book works this way throughout: find a pattern on the keys, hear what it does, then give it a name.

\section*{Where we are going}

\begin{itemize}
  \item \emph{Chapter 1, The keyboard.} Find every note on the piano, name it, and measure the distance between notes.
  \item \emph{Chapter 2, Scales and keys.} Build the major scale from one short pattern, and see what ``the key of G'' means.
  \item \emph{Chapter 3, Chords.} Build major and minor chords by counting keys, and find the chords that belong together.
  \item \emph{Chapter 4, Smooth chord changes.} Rearrange chords so your hand barely moves, and accompany a song with two hands.
  \item \emph{Chapter 5, Rhythm.} Count time: beats, note lengths, rests and time signatures.
  \item \emph{Chapter 6, Reading notes.} Read notes on the staff, the five lines of sheet music.
  \item \emph{Chapter 7, Your first pieces.} Put it all together, first a melody, then both hands, with a short daily practice routine.
\end{itemize}

\section*{Why music is in a mathematics curriculum}

Music is full of patterns: repeating groups of keys, distances you can count, cycles that come back to where they started. You do not need any mathematics to read this book. Counting up to twelve is enough.

If you do enjoy mathematics, look for the boxes titled ``For the mathematically curious''. They explain the numbers behind the sound, such as why some notes blend and others clash. They are optional, and in the web edition they start closed. Skip them freely.
